% 7. Zaključak
\section{Zaključak}
\label{sec:zakljucak}

U ovom radu predstavljen je EduGrader, višemodelska platforma otvorenog koda za ocenjivanje pitanja otvorenog tipa uz podršku veštačke inteligencije i pod nadzorom nastavnika.
Sistem objedinjuje više velikih jezičkih modela (GPT-4o, GPT-4o-mini, DeepSeek-Chat i dva modela usmerena na rezonovanje, GPT-5 i DeepSeek-Reasoner) u podesiv tok ocenjivanja koji podržava evaluaciju i u PRA i u GRA režimu.
EduGrader je evaluiran na 583 autentična studentska odgovora sa pet univerzitetskih kurseva na dve institucije,
koji obuhvataju kratke činjeničke odgovore, teorijska objašnjenja i izraze nalik kodu.

PRA režim daje višu usklađenost sa ljudskim ocenjivačima od GRA režima, kroz sve modele i konfiguracije strogosti.
Modeli usmereni na rezonovanje, naročito GPT-5, najbliži su ljudskom ocenjivanju, sa Pirsonovim koeficijentom korelacije od 0{,}89 u PRA režimu.
Za evaluirane tipove odgovora, a posebno za relativno strukturirana STEM pitanja, savremeni veliki jezički modeli tako aproksimiraju obrasce ljudskog ocenjivanja sa odstupanjima uporedivim sa prirodnom ljudskom varijabilnošću.
Razlike između modela, režima ocenjivanja i konfiguracija strogosti pri tome su sistematske,
pa izbor modela i parametara nije proizvoljan.

Rezultati istovremeno potvrđuju da su sistemi veštačke inteligencije najkorisniji kao pomoćni alat, a ne kao zamena za čoveka.
Ljudski nadzor ostaje neophodan za pravednost, kontekstualno tumačenje i rešavanje graničnih slučajeva.

EduGrader je objavljen kao projekat otvorenog koda, radi reproducibilnosti nalaza i daljeg razvoja sistema.
